
\documentclass[12pt,a4paper]{article}
\usepackage[spanish]{babel}
\usepackage[utf8]{inputenc}
\usepackage[T1]{fontenc}
\usepackage{amsmath, amssymb}
\usepackage{graphicx}
\usepackage{geometry}
\usepackage{cite}
\geometry{margin=2.5cm}

\title{Marco teórico}
\author{}
\date{}

\begin{document}
\maketitle

\section{Introducción general}

La interacción entre la radiación electromagnética y sistemas materiales de tamaño comparable a la longitud de onda constituye un problema fundamental en óptica clásica.
Las partículas esféricas, debido a su simetría geométrica, permiten obtener soluciones analíticas exactas a las ecuaciones de Maxwell, siendo la más importante la
\emph{teoría de Mie}, desarrollada en 1908.
Esta teoría describe rigurosamente los campos internos y dispersados por una esfera homogénea iluminada por una onda plana y constituye el eje central de múltiples áreas
de investigación: óptica atmosférica, fotónica, plasmones, caracterización óptica de materiales y nanoantenas, entre otras.

En este trabajo se estudia el comportamiento de la luz al interactuar con partículas esféricas, apoyándose en dos pilares principales: por un lado, la formulación
electromagnética desarrollada en textos clásicos como \emph{Absorption and Scattering of Light by Small Particles} de Bohren y Huffman, donde se presenta de forma sistemática
la teoría de la dispersión por partículas pequeñas; y por otro lado, los códigos de Fortran desarrollados por Mishchenko y colaboradores, que implementan de manera numéricamente
estable la teoría de Mie para partículas individuales y el método de la matriz--T (T--matrix) para sistemas de múltiples partículas.

El objetivo de este marco teórico es proporcionar una base conceptual sólida que permita comprender:

\begin{itemize}
  \item La física fundamental de la dispersión de la luz por partículas pequeñas.
  \item La formulación electromagnética que conduce a la teoría de Mie para una esfera homogénea.
  \item La definición e interpretación de las secciones eficaces de extinción, absorción y dispersión.
  \item La extensión de la teoría a sistemas formados por múltiples partículas esféricas mediante el método T--matrix.
  \item Los conceptos básicos necesarios para entender el funcionamiento de nanoantenas basadas en partículas esféricas.
  \item La relación entre toda esta teoría y las magnitudes que calcula el código de Fortran de Mishchenko.
\end{itemize}

En las secciones siguientes se desarrollan estos puntos de forma autocontenida, de manera que el lector pueda interpretar correctamente los resultados numéricos obtenidos con
el código de simulación sin necesidad de recurrir constantemente a la bibliografía original.


\section{Fundamentos electromagnéticos de la dispersión}

\subsection{Ecuaciones de Maxwell en medios lineales}

El comportamiento de la radiación electromagnética en un medio lineal, isótropo y homogéneo está gobernado por las ecuaciones de Maxwell en ausencia de cargas y corrientes libres:

\begin{align}
  \nabla \times \mathbf{E} &= -\mu \frac{\partial \mathbf{H}}{\partial t}, \\
  \nabla \times \mathbf{H} &=  \epsilon \frac{\partial \mathbf{E}}{\partial t}, \\
  \nabla \cdot (\epsilon \mathbf{E}) &= 0, \\
  \nabla \cdot (\mu \mathbf{H}) &= 0,
\end{align}

donde $\mathbf{E}$ y $\mathbf{H}$ son los campos eléctrico y magnético, respectivamente, y $\epsilon$ y $\mu$ son la permitividad y permeabilidad del medio.
En medios no magnéticos se suele tomar $\mu \approx \mu_0$, mientras que la respuesta de la materia se incorpora en la permitividad efectiva $\epsilon$.

Suponiendo una dependencia temporal armónica de la forma $\exp(-i\omega t)$, las ecuaciones de Maxwell se reducen a la ecuación de onda vectorial para el campo eléctrico:

\begin{equation}
  \nabla^2 \mathbf{E} + k^2 \mathbf{E} = 0,
\end{equation}

donde

\begin{equation}
  k = \frac{2\pi}{\lambda} n
\end{equation}

es el número de onda en el medio, $\lambda$ la longitud de onda en el vacío y $n$ el índice de refracción del medio circundante.
En materiales absorbentes el índice de refracción es complejo, $m = n + i k$, y la parte imaginaria da cuenta de la absorción interna.

El problema de la dispersión de la luz por una partícula esférica consiste en resolver la ecuación de onda en dos dominios: el interior de la esfera, caracterizado por un índice
de refracción $m$, y el exterior, caracterizado por el índice del medio circundante (típicamente aire o vacío).
Las soluciones deben satisfacer las condiciones de contorno en la superficie de la partícula, que exigen la continuidad de las componentes tangenciales de los campos eléctrico y magnético.

\subsection{Concepto físico de dispersión}

Desde un punto de vista físico, la dispersión puede interpretarse como la re--radiación de energía electromagnética por parte de las cargas ligadas y corrientes de polarización inducidas
en la partícula por el campo incidente.
La onda incidente excita dipolos (y multipolos de orden superior) en el interior de la partícula, que a su vez emiten radiación en todas las direcciones.
En medios disipativos, parte de la energía electromagnética se transforma en energía interna del material, lo que da lugar a la absorción.

El balance energético de la interacción se expresa mediante la relación

\begin{equation}
  C_{\text{ext}} = C_{\text{sca}} + C_{\text{abs}},
\end{equation}

donde $C_{\text{ext}}$ es la sección eficaz de extinción, $C_{\text{sca}}$ la sección eficaz de dispersión y $C_{\text{abs}}$ la de absorción.
Estas magnitudes cuantifican el área efectiva que presenta la partícula frente al haz incidente y constituyen las salidas fundamentales del código de Mishchenko.
El libro de Bohren y Huffman hace especial énfasis en esta interpretación energética, ilustrando con numerosos ejemplos su significado físico en distintos regímenes de tamaño.

\subsection{Parámetro de tamaño y regímenes de dispersión}

Un parámetro adimensional clave en la teoría de la dispersión es el llamado \emph{parámetro de tamaño}

\begin{equation}
  x = k a = \frac{2\pi a}{\lambda},
\end{equation}

donde $a$ es el radio de la partícula.
En función de $x$ se distinguen tres regímenes principales:

\begin{itemize}
  \item $x \ll 1$: régimen de dispersión de Rayleigh, en el que la partícula es mucho más pequeña que la longitud de onda. En este caso la respuesta está dominada por el modo dipolar y las secciones eficaces
        presentan una fuerte dependencia con la longitud de onda ($C_{\text{sca}} \propto \lambda^{-4}$).
  \item $x \sim 1$: régimen de tamaño intermedio, en el que la partícula y la longitud de onda son comparables. En este rango aparecen resonancias asociadas a modos multipolares de orden superior
        y el comportamiento espectral es complejo.
  \item $x \gg 1$: régimen de óptica geométrica, en el que los fenómenos pueden interpretarse en términos de reflexión, refracción y difracción clásicas, y la teoría de Mie se aproxima a los resultados
        de la óptica geométrica.
\end{itemize}

Una de las ventajas de la teoría de Mie es que proporciona una solución exacta válida para cualquier valor de $x$, ofreciendo una descripción continua entre los distintos regímenes.
El código de Fortran de Mishchenko está diseñado para trabajar de forma robusta en un amplio rango de parámetros de tamaño, prestando especial atención a la estabilidad numérica
en el cálculo de funciones especiales de alto orden.



\section{Teoría de Mie para una esfera homogénea}

\subsection{Formulación general y expansión multipolar}

La resolución del problema de dispersión por una esfera homogénea se basa en expandir los campos incidente, interno y dispersado en armónicos esféricos vectoriales.
La simetría esférica del problema permite desacoplar las soluciones en contribuciones eléctricas y magnéticas, que se asocian a modos multipolares de distinto orden.

En la formulación estándar, la onda plana incidente se escribe como una superposición de armónicos esféricos vectoriales regulares, mientras que el campo dispersado se expresa en términos
de armónicos esféricos esféricos salientes (funciones de Hankel).
Aplicando las condiciones de contorno en la superficie de la esfera se obtienen las expresiones analíticas de los coeficientes de Mie $a_n$ y $b_n$, que cuantifican la contribución
de cada modo eléctrico y magnético, respectivamente.

Las expresiones explícitas para estos coeficientes son

\begin{align}
  a_n &= \frac{ m \psi_n(m x)\,\psi_n'(x) - \psi_n(x)\,\psi_n'(m x) }
              { m \psi_n(m x)\,\xi_n'(x) - \xi_n(x)\,\psi_n'(m x) }, \\[4pt]
  b_n &= \frac{ \psi_n(m x)\,\psi_n'(x) - m \psi_n(x)\,\psi_n'(m x) }
              { \psi_n(m x)\,\xi_n'(x) - m \xi_n(x)\,\psi_n'(m x) },
\end{align}

donde $\psi_n$ y $\xi_n$ son las funciones de Riccati--Bessel y la prima indica derivada respecto al argumento.
Estas funciones especiales se relacionan con las funciones de Bessel esféricas de primera especie $j_n$ y con las funciones de Hankel esféricas $h_n^{(1)}$.
En la práctica, el cálculo numérico de $\psi_n$, $\xi_n$ y sus derivadas requiere algoritmos estables para evitar sobreflujos y pérdidas de precisión, especialmente para órdenes altos.

El texto de Bohren y Huffman dedica el capítulo 4 a derivar estas expresiones de manera detallada, discutiendo además diferentes estrategias de truncamiento del sumatorio y criterios
de convergencia.
El código de Mishchenko implementa precisamente estos coeficientes, utilizando algoritmos optimizados para el cálculo recursivo de las funciones de Bessel esféricas, con especial cuidado
en la elección de $n_{\text{max}}$ en función del parámetro de tamaño $x$.

\subsection{Secciones eficaces de dispersión, extinción y absorción}

A partir de los coeficientes de Mie se pueden calcular las principales magnitudes observables que caracterizan la interacción de la luz con la partícula.
Las expresiones para las secciones eficaces de dispersión y extinción son

\begin{align}
  C_{\text{sca}} &= \frac{2\pi}{k^2} \sum_{n=1}^{\infty} (2n+1) \left( |a_n|^2 + |b_n|^2 \right), \\[4pt]
  C_{\text{ext}} &= \frac{2\pi}{k^2} \sum_{n=1}^{\infty} (2n+1) \,\Re(a_n + b_n),
\end{align}

mientras que la sección eficaz de absorción se obtiene como

\begin{equation}
  C_{\text{abs}} = C_{\text{ext}} - C_{\text{sca}}.
\end{equation}

En la práctica resulta útil normalizar estas cantidades dividiéndolas por el área geométrica de la partícula, $\pi a^2$, introduciendo así las \emph{eficiencias ópticas}:

\begin{equation}
  Q_{\text{sca}} = \frac{C_{\text{sca}}}{\pi a^2}, \qquad
  Q_{\text{ext}} = \frac{C_{\text{ext}}}{\pi a^2}, \qquad
  Q_{\text{abs}} = \frac{C_{\text{abs}}}{\pi a^2}.
\end{equation}

El código de Mishchenko devuelve precisamente estas eficiencias como una de sus salidas principales.
De este modo, el usuario puede comparar la respuesta óptica de partículas de distinto tamaño y material de forma adimensional.

\subsection{Matriz de dispersión y parámetros de polarización}

La teoría de Mie permite obtener no solo magnitudes integradas, sino también la distribución angular de la radiación dispersada y su estado de polarización.
Esta información se encapsula en la \emph{matriz de dispersión} o matriz de amplitud, cuyos elementos se suelen denotar por $S_1(\theta)$ y $S_2(\theta)$ para partículas esféricas.
Estos elementos se escriben como

\begin{align}
  S_1(\theta) &= \sum_{n=1}^{\infty} \frac{2n+1}{n(n+1)}
   \left( a_n \pi_n(\theta) + b_n \tau_n(\theta) \right), \\
  S_2(\theta) &= \sum_{n=1}^{\infty} \frac{2n+1}{n(n+1)}
   \left( a_n \tau_n(\theta) + b_n \pi_n(\theta) \right),
\end{align}

donde $\pi_n$ y $\tau_n$ son funciones angulares relacionadas con los polinomios de Legendre asociados.
A partir de $S_1$ y $S_2$ se construyen los elementos de la matriz de Müller, que relaciona los parámetros de Stokes del haz incidente con los del haz dispersado.
Este formalismo es especialmente relevante en problemas donde la polarización juega un papel importante, como en óptica atmosférica o caracterización de aerosoles.

Los programas de Mishchenko calculan los elementos de la matriz de dispersión y de la matriz de Müller para diferentes ángulos de observación, proporcionando así una descripción completa
de la intensidad y el estado de polarización de la luz dispersada.
Esta información es fundamental para comparar con medidas experimentales y para estudiar fenómenos de radiación más complejos.



\section{Interacciones entre múltiples partículas esféricas}

\subsection{Motivación y formulación del problema}

En muchas situaciones físicas de interés no se tiene una única partícula aislada, sino un conjunto de partículas que pueden interactuar entre sí a través de los campos dispersados.
Ejemplos incluyen agregados de nanopartículas, aerosoles densos, suspensiones coloidales o arreglos de nanoantenas.
En estos casos, el campo dispersado por cada partícula actúa como campo incidente sobre las demás, de modo que la solución del problema requiere tener en cuenta la \emph{dispersión múltiple}.

Resolver de forma exacta el problema de muchas partículas es, en general, muy complejo.
Sin embargo, cuando las partículas son esféricas es posible explotar de nuevo la simetría geométrica y utilizar métodos basados en la superposición de soluciones de Mie individuales.
El enfoque más extendido en la literatura es el \emph{método de la matriz--T} (T--matrix method), del cual los códigos de Mishchenko son uno de los referentes principales.

\subsection{Método de superposición T--matrix}

La idea básica del método T--matrix es describir la respuesta de cada partícula mediante una matriz $\mathbf{T}_i$ que relaciona los coeficientes multipolares del campo incidente local
sobre la partícula $i$ con los coeficientes del campo dispersado por ella.
Para una esfera aislada, la matriz $\mathbf{T}_i$ puede construirse directamente a partir de los coeficientes de Mie $a_n$ y $b_n$.
Cuando hay varias partículas, el campo incidente sobre la partícula $i$ es la suma del campo incidente externo y de los campos dispersados por todas las demás partículas:

\begin{equation}
  \mathbf{a}_i^{\text{tot}} = \mathbf{a}_i^{\text{inc}} + \sum_{j \neq i} \mathbf{A}_{ij}\,\mathbf{b}_j,
\end{equation}

donde $\mathbf{a}_i^{\text{tot}}$ son los coeficientes multipolares del campo total incidente sobre la partícula $i$, $\mathbf{b}_j$ los coeficientes del campo dispersado por la partícula $j$
y $\mathbf{A}_{ij}$ son matrices de traslación que permiten expresar los armónicos esféricos centrados en la partícula $j$ en torno al centro de la partícula $i$.

La respuesta de cada partícula viene dada por

\begin{equation}
  \mathbf{b}_i = \mathbf{T}_i \,\mathbf{a}_i^{\text{tot}},
\end{equation}

de modo que sustituyendo se obtiene un sistema lineal acoplado para los vectores $\mathbf{b}_i$:

\begin{equation}
  \mathbf{b}_i = \mathbf{T}_i \left( \mathbf{a}_i^{\text{inc}} + \sum_{j \neq i} \mathbf{A}_{ij}\,\mathbf{b}_j \right).
\end{equation}

En forma matricial, el conjunto de ecuaciones para todas las partículas puede escribirse como

\begin{equation}
  \mathbf{M}\,\mathbf{b} = \mathbf{f},
\end{equation}

donde $\mathbf{M}$ es una matriz densa que contiene las contribuciones de las matrices $\mathbf{T}_i$ y de las matrices de traslación $\mathbf{A}_{ij}$, $\mathbf{b}$ es el vector que agrupa
todos los coeficientes de campos dispersados y $\mathbf{f}$ representa la excitación debida al campo incidente externo.
La resolución de este sistema proporciona los campos dispersados por cada partícula, a partir de los cuales se pueden calcular secciones eficaces globales, funciones de fase,
campos cercanos, etc.

Los códigos de Mishchenko implementan este esquema de forma eficiente y estable, permitiendo simular configuraciones con múltiples partículas esféricas colocadas en posiciones arbitrarias.
Para ello, el programa requiere como entrada las coordenadas $(x_i,y_i,z_i)$ de los centros de cada esfera, sus radios $a_i$ y sus índices de refracción.
A partir de estos datos, el código construye automáticamente las matrices de traslación y ensambla el sistema lineal global.

\subsection{Condición de no solapamiento}

Un aspecto importante de la formulación es que la teoría supone que las partículas no se solapan geométricamente; es decir, la distancia entre los centros de dos esferas cualesquiera
debe satisfacer

\begin{equation}
  |\mathbf{r}_i - \mathbf{r}_j| > a_i + a_j.
\end{equation}

Si esta condición no se cumple, las superficies de las esferas se intersecan y el modelo pierde consistencia física, ya que las condiciones de contorno se vuelven ambiguas.
Por este motivo, los códigos de Mishchenko incluyen comprobaciones explícitas de solapamiento: si se detecta que alguna pareja de partículas viola la desigualdad anterior,
el programa aborta la ejecución y devuelve un mensaje de error.

En el contexto de este trabajo, esta restricción se traduce en una condición geométrica que debe respetarse al generar las configuraciones de partículas que se utilizan como entrada
del código de Fortran.



\section{Bases teóricas para nanoantenas esféricas}

\subsection{Resonancias multipolares en partículas esféricas}

Las partículas esféricas de tamaño comparable a la longitud de onda pueden presentar resonancias asociadas a los diferentes modos multipolares descritos por los coeficientes de Mie.
En el régimen $x \sim 1$ la respuesta óptica está dominada por un conjunto discreto de modos resonantes:

\begin{itemize}
  \item El modo dipolar eléctrico, asociado principalmente al coeficiente $a_1$.
  \item El modo dipolar magnético, asociado al coeficiente $b_1$, especialmente relevante en materiales dieléctricos de alto índice.
  \item Modos de orden superior (cuadrupolos, octupolos, etc.), que corresponden a $a_n$ y $b_n$ con $n>1$.
\end{itemize}

En materiales metálicos, la resonancia dipolar eléctrica está íntimamente relacionada con los plasmones de superficie localizados, mientras que en materiales dieléctricos de alto índice
(Si, Ge, TiO$_2$, etc.) destacan las resonancias magnéticas.
Estas resonancias producen fuertes variaciones en las secciones eficaces y en los campos cercano y lejano, y constituyen el mecanismo físico básico de funcionamiento de las nanoantenas ópticas.

\subsection{Intensificación del campo cercano}

La teoría de Mie permite calcular no solo el campo lejano, responsable de la radiación dispersada medida a grandes distancias, sino también el campo electromagnético en las proximidades
inmediatas de la partícula.
En las cercanías de una resonancia, el campo eléctrico puede verse fuertemente intensificado en regiones concretas alrededor de la partícula, dando lugar a efectos como la concentración
de energía o el aumento de la densidad local de estados electromagnéticos.

Los enfoques basados en la matriz--T, como el de Mishchenko, permiten evaluar estos campos cercanos en puntos arbitrarios del espacio, tanto para una única partícula como para
conjuntos de partículas acopladas.
Esta capacidad es fundamental para el diseño de nanoantenas, ya que el rendimiento de muchas aplicaciones (emisión dirigida, detección, sensado) depende directamente
del perfil espacial del campo cercano.

\subsection{Direccionalidad y patrones de radiación}

Otra propiedad relevante de las nanoantenas es su capacidad para modificar la distribución angular de la radiación, concentrando la emisión en determinadas direcciones
o suprimiéndola en otras.
En el caso de partículas esféricas, la direccionalidad se puede analizar a partir de los elementos $S_1(\theta)$ y $S_2(\theta)$ de la matriz de dispersión.
Determinadas combinaciones de los modos eléctricos y magnéticos pueden dar lugar a patrones de radiación fuertemente asimétricos, incluyendo la supresión de la retrodispersión
o de la dispersión hacia adelante (condiciones de Kerker).

Por ejemplo, cuando el modo dipolar eléctrico y el modo dipolar magnético están en fase y de amplitud comparable ($a_1 \approx b_1$) se puede obtener una dispersión predominantemente
hacia adelante, mientras que si están en oposición de fase se favorece la dispersión hacia atrás.
Estas ideas se utilizan extensamente en el diseño de nanoantenas dieléctricas para lograr direccionalidad en la emisión o en la recolección de luz.

Los códigos basados en la matriz--T permiten calcular de forma directa los patrones de radiación de sistemas de partículas esféricas, de modo que es posible estudiar configuraciones
que actúan como nanoantenas simples o arreglos más complejos (dímeros, cadenas, etc.).
En el contexto de este trabajo, la teoría expuesta proporciona la base conceptual necesaria para interpretar los resultados de dispersión y campos cercanos
que posteriormente se analizarán en configuraciones con potencial aplicación como nanoantenas.


\section{Relación entre la teoría y el código de Fortran de Mishchenko}

\subsection{Parámetros de entrada y su significado físico}

El código de Fortran de Mishchenko requiere como entrada un conjunto de parámetros que se corresponden directamente con las magnitudes introducidas a lo largo de esta sección teórica.
Entre los más relevantes se encuentran:

\begin{itemize}
  \item El número de partículas esféricas presentes en la configuración.
  \item El radio $a_i$ de cada esfera y su índice de refracción complejo $m_i = n_i + i k_i$ a la longitud de onda de interés.
  \item La longitud de onda $\lambda$ (o, de forma equivalente, la frecuencia o el número de onda $k$).
  \item Las posiciones de los centros de las partículas $(x_i, y_i, z_i)$ en el sistema de coordenadas global.
  \item La dirección de propagación y el estado de polarización de la onda incidente.
  \item El orden de truncamiento máximo $n_{\text{max}}$ para las expansiones multipolares.
\end{itemize}

Estos parámetros permiten definir por completo el problema físico de dispersión, de acuerdo con la teoría de Mie y el método T--matrix.
El usuario debe elegirlos de manera coherente con la configuración que se desea estudiar y con las condiciones experimentales que se pretenden modelizar.

\subsection{Cálculos internos realizados por el código}

Una vez proporcionados los parámetros de entrada, el código de Mishchenko realiza internamente una serie de pasos que pueden resumirse del siguiente modo:

\begin{enumerate}
  \item Para cada tipo de partícula, calcula los coeficientes de Mie $a_n$ y $b_n$ a partir del radio, el índice de refracción y la longitud de onda.
  \item Construye la matriz--T individual de cada esfera, utilizando la relación entre los coeficientes de Mie y los coeficientes multipolares de los campos incidente y dispersado.
  \item Evalúa las matrices de traslación $\mathbf{A}_{ij}$ necesarias para expresar los campos dispersados por una partícula como campos incidentes sobre las demás.
  \item Ensambla el sistema lineal global que describe la dispersión múltiple en el conjunto de partículas y lo resuelve mediante técnicas numéricas apropiadas
        (factorización directa o métodos iterativos, en función del tamaño del problema).
  \item A partir de los coeficientes de campo dispersado $\mathbf{b}_i$ calcula las secciones eficaces globales del sistema, la matriz de dispersión para diferentes ángulos,
        y, si se desea, los campos cercanos en puntos específicos del espacio.
\end{enumerate}

Cada uno de estos pasos tiene una base teórica directa en los desarrollos presentados en las secciones anteriores.
Por tanto, el marco teórico proporciona la interpretación física de las magnitudes que el código calcula de manera numérica.

\subsection{Salidas típicas y su interpretación}

Entre las salidas más habituales del código se encuentran:

\begin{itemize}
  \item Las eficiencias ópticas globales del sistema ($Q_{\text{ext}}$, $Q_{\text{sca}}$, $Q_{\text{abs}}$), que permiten cuantificar el balance energético de la interacción luz--partículas.
  \item La matriz de Müller o matriz de dispersión para diferentes ángulos de observación, a partir de la cual se obtienen patrones de radiación e información sobre la polarización.
  \item Los campos cercanos en mallas bidimensionales o tridimensionales alrededor de las partículas, que permiten visualizar intensificaciones locales del campo
        y modos resonantes relevantes para aplicaciones de nanoantenas.
\end{itemize}

La comparación entre estos resultados numéricos y las predicciones analíticas de la teoría de Mie, así como con observaciones experimentales, posibilita la validación del modelo
y la exploración de nuevos fenómenos físicos en sistemas de partículas esféricas.

En resumen, el código de Fortran de Mishchenko constituye una implementación computacional directa de la teoría electromagnética desarrollada en este marco teórico.
El conocimiento de dicha teoría es imprescindible para seleccionar adecuadamente los parámetros de entrada, interpretar correctamente las salidas del programa
y diseñar configuraciones de partículas con propiedades ópticas específicas, como las que se buscan en el contexto de nanoantenas y otras aplicaciones fotónicas.

\end{document}
